\section{Background and Motivation}
\label{sec:background}

\subsection{From function invocation to sandbox session}

A function-as-a-service platform offers one primitive: invoke a function,
return a result. Its benchmark suites follow that shape. SeBS measures
invocation latency, cold-start cost and resource use across six application
classes~\cite{sebs2021}; FunctionBench provides micro and application
workloads for the same model~\cite{functionbench2019}; vHive and the snapshot
line of work attack the cold start directly~\cite{vhive2021,catalyzer2020}.

A sandbox platform offers five: create, exec, suspend, resume, destroy. The
execution model is
\[
  \text{create} \rightarrow (\text{exec})^{n} \rightarrow
  [\text{suspend} \rightarrow \text{resume} \rightarrow (\text{exec})^{m}]
  \rightarrow \text{destroy},
\]
with $n$ and $m$ in the hundreds for an agent session and the bracketed
segment repeated for a long-lived one. Three measurement consequences follow.

First, cold start stops being the headline. It is paid once per session and
amortised over $n$ commands, so for large $n$ the per-command overhead
dominates the user-visible total. A ranking built on create latency alone can
inverse the ranking built on session completion time. \needsdata{whether it
does in practice --- this is RQ1}

Second, the exec path becomes a first-class object of study. It is not part of
the FaaS model at all, and providers have chosen materially different
mechanisms for it (\S\ref{sec:implementation}).

Third, correctness joins the measurement. Suspend and resume make the
filesystem a thing that can be silently lost, so a lifecycle benchmark that
reports only latency cannot distinguish a fast resume from a fresh boot.

\subsection{What the workloads actually do}
\label{sec:background:characterisation}

Benchmark design without workload data is guesswork dressed as rigour. Before
fixing the suite we therefore need answers to:

\begin{enumerate}
  \item How many commands does one sandbox session run, and how is that
  distributed?
  \item How are individual command durations distributed? Where is the tail?
  \item What fraction of session wall-clock goes to environment preparation,
  to useful execution, and to waiting for the next interaction?
  \item How often is a reinforcement-learning environment reset, and by which
  mechanism?
  \item How large is the state that must survive a suspend?
  \item What is the concurrency and the request arrival pattern of a real
  rollout fleet?
\end{enumerate}

\needsdata{answers from agent execution traces, public RL environments, and
aggregated provider metrics. Where a trace cannot be obtained, a synthetic
workload is used and labelled as synthetic, with its representativeness stated
as a limitation rather than assumed.}
